% =====================================================================
% Spin III: The ZBW Mode Spectrum and the Lattice Voronoi-Cell
% Eigenvalue Problem
% Third paper of the CPP spin sub-arc (Spin I, Spin II, Spin III).
%
% CHANGELOG
%   V0   19 Aug 2026  Patch 3234  Initial draft (Session 149), written
%        from the arc's own materials at founder request: the
%        development-notes charter ("Spin III: derive the sub-harmonic
%        spectrum from the 600-cell Voronoi eigenvalues"), the committed
%        eigenvalue script (16 Mar 2026), its data and figure. No Spin
%        III paper previously existed (full-history finding, Patch 3232
%        record). Honest-content ruleset: the rigorous pieces (mode
%        arithmetic, FD verification, scale-bridge identity,
%        discreteness bound) are stated as results; the lattice-cell
%        capstone is reported as OPEN with the committed instrument's
%        boundary-condition mismatch diagnosed and a corrected
%        instrument frozen (OPEN-QM-8). Verify:
%        scripts/3234_verify_spin3_v0.py (10/10).
%   V0.1 19 Aug 2026  Patch 3236  A1 RULED (founder: true Voronoi cell;
%        24-cell was a prior worker's initiative) and the corrected
%        measurement EXECUTED: sphere control exact; true cell
%        (regular dodecahedron, dual-120-cell) returns the frozen
%        verdict MODE2-RECOVERED at both densities (node 0.6670 vs
%        2/3; antinode 0.3333; spectrum shift 0.31%). New Section 7;
%        ledger updated. Run: scripts/3236_qm8_true_cell_run.py;
%        record: qm8_corrected_run_record.md.
%   V0.2 19 Aug 2026  Patch 3238  THE ANALYTIC LEG (founder go):
%        Section 8 — icosahedral selection and protection. Exact
%        character counts (m_l = 0 for l = 1..5; first invariant
%        channel l = 6); the cell's measured harmonic content
%        (eps_6 = 0.051, forbidden channels at quadrature floor);
%        the Selection Theorem (k(6,1) = 8.211 > 3pi/2, so the
%        second invariant mode IS Mode 2, with the eight global
%        interlopers identified as the l=1+l=2 multiplets matching
%        the 3236 indices); the Protection Theorem (positions shift
%        at O(eps_6^2); predicted 0.0026 vs measured 0.0031).
%        Verify: scripts/3238_verify_2I_selection.py (7/7).
%   V0.3 19 Aug 2026  Patch 3244  CONV-026 calibrations (Q5 4-1
%        VALID; Q6 3-2 CALIBRATED; Q8 3-2 RATIFY — status
%        panel-ratified): worked character table printed in 8.1
%        (Copilot); exact Bessel-route gap 8.2108 recorded beside FD
%        (GPT; three routes, one ordering); Y_l orthogonality line in
%        the Protection proof (DeepSeek mini-lemma); DeepSeek's
%        stricter consistency sentence + the residual's specified
%        discharge path in 8.4. No claim strengthened; two
%        residual notes now on OPEN-QM-8 (non-perturbative bound
%        with discharge path; dynamic symmetry-breaking robustness,
%        Gemini provenance).
%   v1.0.1 19 Aug 2026 Patch 3253  RE-IDENTIFIED: permanent ID SPIN-3
%        (founder naming confirmation, Session 150); folder/file renamed
%        spin-III_600cell_voronoi_ZBW_eigenvalues -> SPIN-3_...; title ID
%        prefixed; title version line corrected (had lagged at 'Version 0
%        (draft)' after the v1.0 ship - bookkeeping defect caught here);
%        self-citations tagged with SPIN-1/SPIN-2. Content unchanged.
%   v1.0 19 Aug 2026  Patch 3248  SHIPPED. Paper production protocol
%        executed: Keywords + Plain Language Summary + CPP Physical
%        Mechanism Bridge & conventional-mapping section added
%        (Phase 2 formatting); 8-file documentation suite produced
%        (Phase 7); Phase 5 N/A (no new axioms/predictions); Phase 6
%        deposit = founder/Isak Zenodo test run, BLOCKED on the
%        SPIN-N naming confirmation + DOI reservation; Phase 8
%        transcript curation registered as owed. Review basis:
%        CONV-026 (Q5 4-1 VALID; Q8 3-2 RATIFY; cycle closed 3247).
% =====================================================================
\documentclass[12pt]{article}

\usepackage{amsmath,amssymb,amsthm}
\usepackage{geometry}
\usepackage{hyperref}
\usepackage{booktabs}
\usepackage{enumitem}
\usepackage{parskip}
\geometry{margin=1in}

% ---- corpus macros (matching the Spin I/II conventions) ----
\newcommand{\lP}{l_P}
\newcommand{\SSV}{\text{SSV}}
\newcommand{\rth}{r_{\rm th}}
\newcommand{\rin}{r_{\rm in}}
\newcommand{\rout}{r_{\rm out}}
\newcommand{\rC}{r_C}
\newcommand{\aBohr}{a_0}

\newtheorem{theorem}{Theorem}
\newtheorem{proposition}{Proposition}
\newtheorem{lemma}{Lemma}
\newtheorem{remark}{Remark}

\title{SPIN-3: The ZBW Mode Spectrum and the Lattice
Voronoi-Cell Eigenvalue Problem\\
\normalsize Version 1.0.1 --- Conscious Point Physics spin sub-arc (Spin III)}
\author{Thomas Lee Abshier, ND\\ Hyperphysics Institute}
\date{19 August 2026}

\begin{document}
\maketitle

\begin{abstract}
This paper is the capstone of the CPP spin sub-arc.  Spin~I established
that the captured Dipole Particle orbital geometry yields the electron
spin angular momentum $L = \hbar/2$ exactly, \emph{given} the radius
condition $\rout = 2\rin$.  Spin~II derived that condition from the ZBW
polarization cloud's standing-wave structure: an open--closed resonator
(node at $r=0$ by CP Exclusion, antinode at the thermal boundary
$\rth = \hbar/2m_ec$) whose Mode~2 has an interior antinode at $\rth/3$
and an interior node at $2\rth/3$ --- ratio $2$, exactly, independent of
every physical constant.  Spin~III's charter, set when Spin~II was
written, is to ground that mode structure in the substrate itself: to
derive the spectrum from the eigenvalue problem of the lattice
Voronoi cell.  This V0 establishes what can currently be established
and reports plainly what cannot.  \textbf{Established here:} the exact
resonator spectrum $k_n = (2n-1)\pi/2\rth$ with finite-difference
verification to $0.01\%$ across eight modes by two independent
constructions; the Mode-2 geometry to $0.1\%$; the scale-bridge
identity $\rin^{\rm (I)}/\rin^{\rm (II)} = 6/[4\alpha(1+\sqrt{2})^2]
= 35.2675$, shown to be an \emph{exact algebraic identity} rather than
a numerical coincidence; and the discreteness bound
$(\lP/\rth)^2 \approx 7\times10^{-45}$, which fixes the division of
labor --- the lattice must supply mode \emph{selection}, not mode
\emph{corrections}.  \textbf{Diagnosed here:} the arc's committed
lattice-cell computation (a graph Laplacian on the 24-cell interior)
solves the closed (Neumann-type) problem rather than the open--closed
one --- the CP-Exclusion node at $r=0$ is never encoded --- so its low
spectrum consists of the four 4D dipole modes and its own Mode-2
diagnostic correctly reports no identification.  \textbf{Measured here
(V0.1):} the founder ruled the domain question --- the 24-cell carried
no physical-picture weight; the true Voronoi cell of the 600-cell-based
lattice is the \emph{regular dodecahedron} (dual-120-cell cell) --- and
the corrected instrument, frozen before its data existed, was executed:
the sphere control reproduces the continuum result to five decimal
places, and on the true cell the frozen verdict fires
\textsc{mode2-recovered} at both mesh densities (Mode-2 node at
$0.6670$ against the exact $2/3$; antinode at $0.3333$; spectrum
shifted $0.31\%$ from the sphere).  The lattice cell's icosahedral
anisotropy does not disturb the mode topology the captured DP anchors
to.  \textbf{Proved here (V0.2), the analytic leg:} the icosahedral
rotation group admits \emph{no} invariant spherical harmonic for
$1 \le l \le 5$ (exact character counts), so the cell's anisotropy
lives entirely in $l = 6, 10, 12, \ldots$ (measured on the actual
cell: $\epsilon_6 = 0.051$, forbidden channels at the quadrature
floor); within the invariant sector the first anisotropy-carrying mode
sits at $kR = 8.211$, far above Mode 2's $3\pi/2 = 4.712$, so
\emph{the second invariant mode of the true cell is the anchoring
mode --- the lattice must choose Mode 2} (Selection Theorem); and the
radial family's fractional node/antinode positions are protected to
$O(\epsilon_6^2)$ (Protection Theorem), with the predicted
$\epsilon_6^2 = 0.0026$ matching the measured $0.0031$ eigenvalue
residual to $20\%$.  What remains registered is only a
non-perturbative-in-$\epsilon_6$ bound --- a rigor upgrade, not a gap
in the physical argument.
\end{abstract}

\noindent\textbf{Keywords:} Zitterbewegung; electron spin; 600-cell
lattice; Voronoi cell; regular dodecahedron; icosahedral symmetry;
mode selection; boundary perturbation; Conscious Point Physics.

\section*{Plain Language Summary}
The electron's spin, in this framework, comes from a tiny standing wave
--- like the vibration of an organ pipe closed at one end --- inside the
electron's jitter cloud.  The first two papers of this arc derived the
spin value and the wave's shape assuming space is smooth.  But in CPP,
space is a lattice of discrete points, and each point's territory is a
twelve-sided room (a regular dodecahedron), not a perfect sphere.  This
paper asks: does the wave pattern survive in the real room?  We answer
three ways.  First, the room's corners are astronomically too small to
distort the wave (forty-plus orders of magnitude).  Second, we build a
careful computational instrument, test it on a perfect sphere where the
answer is known exactly, then run it in the twelve-sided room: the
pattern survives, its landmarks moving by less than a tenth of a
percent.  Third --- the elegant half --- we prove with symmetry
arithmetic that the room \emph{cannot} support any competing pattern:
the icosahedral shape forbids every lopsided wave below a high
threshold, so the lattice has no choice but to select exactly the wave
that produces spin-$\tfrac12$.  The proof even predicted the size of the
tiny imperfections the computer had already measured.

\tableofcontents

% =====================================================================
\section{The arc and this paper's charter}
\label{sec:arc}

The spin sub-arc was planned as three papers, and the plan is recorded
verbatim in the arc's development notes:
\begin{itemize}[leftmargin=2em]
\item \textbf{Spin I}: given $\rout = 2\rin$, derive $L = \hbar/2$.
The captured-DP orbital geometry with the $1/r^2$ force law gives the
angular-frequency ratio $\omega_{\rm in}/\omega_{\rm out} = 2\sqrt{2}$
and total orbital angular momentum $\hbar/2$ exactly when
$\rin = \aBohr/[4(1+\sqrt{2})^2]$.
\item \textbf{Spin II}: derive $\rout = 2\rin$ from ZBW sub-harmonics.
The ZBW cloud is an open--closed radial resonator; its Mode~2 anchors
the inner $+$CP at the interior antinode ($\rth/3$) and the outer $-$CP
at the interior node ($2\rth/3$).
\item \textbf{Spin III}: derive the sub-harmonic spectrum from the
lattice Voronoi-cell eigenvalues --- ``the last mathematical gap to the
lattice.''
\end{itemize}
The honesty constraint on this third paper was also set in advance, in
the same notes: the lattice claim must not be asserted without either a
numerical eigenvalue computation of the lattice cell or an analytic
selection argument from the $2I$ symmetry group.  This V0 obeys that
constraint.  It consolidates what is provable now, diagnoses why the
existing numerical attempt did not settle the capstone, and freezes the
corrected instrument that can.

% =====================================================================
\section{The continuum spectrum, verified two ways}
\label{sec:spectrum}

\begin{proposition}[Resonator spectrum; Spin II]
The ZBW radial problem in the $\ell = 0$ sector, with a node at $r = 0$
(CP Exclusion) and a free antinode at $r = \rth$, has modes
\begin{equation}
\psi_n(r) = \sin(k_n r), \qquad
k_n = \frac{(2n-1)\pi}{2\rth}, \quad n = 1, 2, 3, \ldots
\label{eq:spectrum}
\end{equation}
\end{proposition}

\begin{theorem}[Mode-2 geometry; Spin II]
Mode 2 has its interior antinode at $\rin = \rth/3$ and its interior
node at $\rout = 2\rth/3$, whence $\rout/\rin = 2$ exactly ---
integer arithmetic of trigonometric zeros, independent of every
physical constant.
\end{theorem}

Table~\ref{tab:fd} reproduces the arc's committed finite-difference
verification (forward-difference Neumann ghost; the committed CSV was
regenerated bit-identically this session) alongside this paper's
independent reconstruction (symmetric lumped-mass Neumann formulation
--- a different discretization route).  Both agree with
Eq.~\eqref{eq:spectrum} at the $10^{-4}$ level across eight modes, and
the independent route locates the Mode-2 antinode and node at $R/3$ and
$2R/3$ to $0.1\%$ with ratio $2.000$ (verify script, checks 1--2).

\begin{table}[h]
\centering
\caption{Resonator eigenvalues: analytic vs.\ committed FD
(\texttt{data/spin3\_eigenvalues.csv}).  The independent lumped-mass
reconstruction agrees with both to $<0.05\%$
(\texttt{scripts/3234\_verify\_spin3\_v0.py}).}
\label{tab:fd}
\begin{tabular}{cccc}
\toprule
mode $n$ & $k_n$ (analytic) & $k_n$ (FD) & rel.\ error \\
\midrule
1 & 1.57079633 & 1.57063926 & $1.0\times10^{-4}$ \\
2 & 4.71238898 & 4.71191761 & $1.0\times10^{-4}$ \\
3 & 7.85398163 & 7.85319551 & $1.0\times10^{-4}$ \\
4 & 10.99557429 & 10.99447263 & $1.0\times10^{-4}$ \\
5 & 14.13716694 & 14.13574866 & $1.0\times10^{-4}$ \\
6 & 17.27875959 & 17.27702330 & $1.0\times10^{-4}$ \\
7 & 20.42035225 & 20.41829623 & $1.0\times10^{-4}$ \\
8 & 23.56194490 & 23.55956715 & $1.0\times10^{-4}$ \\
\bottomrule
\end{tabular}
\end{table}

% =====================================================================
\section{The scale-bridge identity}
\label{sec:bridge}

Spin~I places the inner orbit at the Bohr-derived scale
$\rin^{\rm (I)} = \aBohr/[4(1+\sqrt{2})^2]$; Spin~II places it at the
ZBW-cloud scale $\rin^{\rm (II)} = \rth/3$.  The committed script
records their ratio as $35.267491$ and observes that it equals
$6/[\alpha\cdot4(1+\sqrt{2})^2]$.  This paper upgrades that observation
to a theorem-grade statement:

\begin{theorem}[Scale bridge, exact]
Given the standard relations $\aBohr = \rC/\alpha$ and $\rth = \rC/2$,
\begin{equation}
\frac{\rin^{\rm (I)}}{\rin^{\rm (II)}}
= \frac{\aBohr/[4(1+\sqrt{2})^2]}{\rth/3}
= \frac{3\,\aBohr}{4(1+\sqrt{2})^2\,\rth}
= \frac{6}{4\alpha(1+\sqrt{2})^2}
\approx 35.2675
\label{eq:bridge}
\end{equation}
exactly --- an algebraic identity, not a numerical coincidence.
\end{theorem}

\noindent (Verified numerically to $10^{-12}$; verify check~3.)  The
physical reading, following the arc's development record: the
sub-harmonic constrains the \emph{ratio}, the Coulomb force balance
constrains the \emph{scale}, and the bridge between the two scales
involves $\alpha$ in exactly the combination \eqref{eq:bridge}.  The
mode structure and the orbital scale are two constraints, not one
constraint stated twice, and their compatibility is what
Eq.~\eqref{eq:bridge} records.

% =====================================================================
\section{The division of labor: selection, not corrections}
\label{sec:division}

The discreteness parameter at the ZBW scale is
\begin{equation}
\left(\frac{\lP}{\rth}\right)^2 \approx 7.0\times10^{-45},
\end{equation}
(verify check~3), so lattice discreteness cannot measurably shift the
continuum eigenvalues \eqref{eq:spectrum} at any accessible precision.
This sharpens what Spin~III must actually deliver: the lattice's job is
not to \emph{correct} the spectrum but to \emph{select} it --- to show
that the substrate's cell geometry, with the CP-Exclusion node and the
thermal free boundary as its physical boundary conditions, produces the
open--closed mode family (and in particular admits the Mode-2
node/antinode topology the captured DP anchors to) rather than some
other family.  A selection statement survives at any lattice fineness;
a correction statement would be empty at $10^{-45}$.

% =====================================================================
\section{The committed lattice-cell computation: what it did and did
not show}
\label{sec:instrument}

The arc's committed script (16 March 2026) contains, as its Part~3, a
lattice-cell computation: rejection-sample the interior of the 24-cell
(unit circumradius, the origin included as a sample point), build a
Gaussian-kernel $k$-nearest-neighbour graph, take the normalised graph
Laplacian $L_{\rm sym} = I - D^{-1/2}WD^{-1/2}$, and inspect the
radially averaged low eigenmodes for a ``Mode-2-like'' profile (exactly
one interior radial sign change).  Executed this session at both
reduced and default sampling ($8\,000$ and $20\,000$ points), the
computation reports, by its own diagnostic, \emph{no Mode-2-like mode}
--- and this V0's finding is that it cannot, for a stated reason:

\begin{proposition}[Boundary-condition mismatch]
The kNN graph Laplacian on uniformly sampled interior points
approximates the Laplacian with \emph{closed} (Neumann-type, free)
boundary behaviour everywhere; nothing in the construction encodes the
CP-Exclusion Dirichlet node at $r = 0$.  Its lowest nontrivial modes
are therefore the four 4D dipole modes of the closed cell, not the
open--closed radial family.
\end{proposition}

\noindent The diagnosis is verified reproducibly (verify check~4, same
construction, same seed): the lowest nontrivial spectrum is a 4-fold
near-degenerate group (spread $\sim10\%$, gap ratio $\sim2$ to the next
group --- the dipole signature in four dimensions), and no mode among
the lowest ten has exactly one interior radial zero.  Two consequences,
stated plainly:
\begin{itemize}[leftmargin=2em]
\item The Part-3 computation, as committed, neither confirms nor
disconfirms the capstone claim.  It solves a different boundary-value
problem.  (The committed figure, consistently, carries panels for
Parts~1, 2 and 4 only.)
\item Nothing upstream is affected: Parts~1--2 (the continuum spectrum
and its FD verification) are independent of Part~3 and reproduce
bit-identically.
\end{itemize}

\subsection{The domain question (founder ruling requested)}
\label{sec:domain}

The committed computation takes the \textbf{24-cell} as its domain.
The arc's charter language says ``600-cell Voronoi eigenvalue
spectrum.''  These are different objects, and V0 declines to paper over
the difference: whether the 24-cell is (a)~the intended Voronoi domain
of the CPP Grid-Point lattice, (b)~a deliberate tractable surrogate for
it, or (c)~a stand-in that should be replaced by the true GP-lattice
Voronoi cell, is a physical-picture question reserved to the founder.
\textbf{Ruled (19 Aug 2026, registered
\texttt{founders\_voice/founder\_ruling\_A1\_voronoi\_domain\_2026-08-19.md}):}
the 24-cell was a prior worker's initiative with no physical-picture
weight; the corrected run uses the true Voronoi cell.  The geometry is
exact and standard: the 600-cell's dual is the 120-cell, so the Voronoi
domain of each lattice vertex is the corresponding \emph{regular
dodecahedron} --- twelve pentagonal faces whose normals point at the
twelve icosahedral nearest-neighbor directions (the vertex figure of
FACT G1).

% =====================================================================
\section{The corrected instrument, frozen before its data exists}
\label{sec:corrected}

Per corpus practice (freeze the repair by principle, before the testing
data exists), the corrected lattice-cell measurement is specified now:

\begin{enumerate}[leftmargin=2em]
\item \textbf{Boundary conditions encoded, not implied.}  Dirichlet on
a small central core (all sample points with $r < \epsilon$ pinned to
zero --- the CP-Exclusion node), free (natural/Neumann) outer boundary
--- the open--closed problem the physics states.
\item \textbf{Radial-mode identification by isotropy score, not raw
sign counting.}  For each low eigenmode, compute the fraction of its
variance explained by its radial average; score modes with fraction
above a declared threshold as radial candidates; count interior zeros
on those only.  (The committed sign-counter applied to raw radial
averages of anisotropic modes is what produced zero counts of 13--36.)
\item \textbf{Domain per the founder's A1 ruling}: 24-cell, or the
GP-lattice Voronoi cell, or both as a robustness pair.
\item \textbf{Frozen readings.}  \textsc{mode2-recovered}: a radial
candidate with exactly one interior zero, node in $[0.60, 0.73]$ and
antinode in $[0.28, 0.40]$ of the boundary radius, stable across two
seeds and two sampling densities.  \textsc{mode2-absent}: no such mode
among the lowest twenty radial candidates at the highest density.
\textsc{indeterminate}: anything else $\Rightarrow$ instrument work,
not a verdict.
\item \textbf{Worker expectation, declared pre-run
(anti-extraction).}  \textsc{mode2-recovered} --- the continuum limit
is exact, discreteness corrections sit at $10^{-45}$, and the closed
problem's dipole modes are removed by the Dirichlet core; the risk is
angular interleaving in four dimensions, which the isotropy score is
designed to separate.
\end{enumerate}

This measurement, plus the analytic $2I$-symmetry selection argument
(the second route the arc's notes named), constitute
\texttt{OPEN-QM-8}.  Neither is attempted in V0.

% =====================================================================
\section{The corrected measurement, executed (V0.1)}
\label{sec:measurement}

The instrument of Section~\ref{sec:corrected} was implemented on the
$u$-equation directly, $-u'' - r^{-2}\Delta_S u = k^2 u$, with the
Dirichlet center and the free end $u'=0$ along each ray --- closing not
only the missing-node defect of Section~\ref{sec:instrument} but a
second semantic trap recorded here for the reviewer: a free (Neumann)
condition on $\psi$ is \emph{not} Spin II's free condition on $u$ (on
a sphere, $\psi$-Neumann gives the $\tan kR = kR$ spectrum,
$k_1R \approx 4.493$, not $(2n-1)\pi/2$).  Discretization: cotangent
Laplacian on an icosphere for $\Delta_S$; per-ray FEM in the
fractional coordinate $s = r/R(\omega)$; boundary radius by the
dodecahedron's support function over the twelve icosahedral face
normals.  Declared approximations (anisotropy-squared class,
$\sim1\%$ in positions): shell-averaged angular metric weight;
ray-coordinate cross-terms neglected.

\textbf{Sphere control (instrument validation):} $k_1 R = 1.57079$
($\pi/2 = 1.57080$), $k_2 R = 4.71209$ ($3\pi/2$), Mode-2 zero at
$0.66667$, antinode at $0.3333$ --- exact to instrument precision.

\textbf{True Voronoi cell, two densities ($162\times80$ and
$642\times120$):} three radial candidates among the lowest ninety
modes at both densities (isotropy scores $0.999$--$1.000$); Mode-2
$k\langle R\rangle = 4.6976$ ($-0.31\%$ from the sphere); node at
$s = 0.6670$ against the exact $2/3$; antinode at $0.3333$ against
$1/3$; Mode~3 exhibiting the expected double-zero structure
($0.402, 0.801$ against $0.4, 0.8$).  \textbf{Frozen verdict
(Section~\ref{sec:corrected} readings, applied verbatim):
\textsc{mode2-recovered}} at both densities, node drift $0.0000$
against the $0.02$ gate.  The worker's expectation, declared at V0
before the A1 ruling existed, is confirmed --- and, more to the point,
the icosahedral cell's $\pm11\%$ radial anisotropy moves the Mode-2
anchoring geometry by less than $0.1\%$: the lattice cell
\emph{selects} the open--closed family and preserves exactly the
topology the captured DP anchors to.  Full record:
\texttt{qm8\_corrected\_run\_record.md}; instrument:
\texttt{scripts/3236\_qm8\_true\_cell\_run.py} (5/5 checks).

The analytic leg --- the $2I$-symmetry selection argument the arc's
notes demanded --- is Section~\ref{sec:analytic}.

% =====================================================================
\section{The analytic leg: icosahedral selection and protection (V0.2)}
\label{sec:analytic}

This section supplies the symmetry argument the arc's development notes
named as Spin III's remaining obligation: not merely that the true cell
\emph{happens} to preserve the Mode-2 structure (the measurement of
Section~\ref{sec:measurement}), but that the lattice's icosahedral
symmetry \emph{forces} the selection.  Every arithmetic claim below is
a checked computation in
\texttt{scripts/3238\_verify\_2I\_selection.py} (7/7).

\subsection{Lemma: the invariant content of the icosahedral group}
\label{sec:characters}

Let $m_l$ be the multiplicity of the trivial representation of the
icosahedral rotation group $I$ (order 60) in the spin-$l$
representation of $SO(3)$:
\begin{equation}
m_l = \frac{1}{60}\sum_{\text{classes}} |C|\,
\frac{\sin\!\big((l+\tfrac12)\theta_C\big)}{\sin(\theta_C/2)},
\end{equation}
over the classes $\{E, 12\,C_5, 12\,C_5^2, 20\,C_3, 15\,C_2\}$.  Exact
evaluation gives
\begin{equation}
m_l = 0 \ \text{for } 1 \le l \le 5;\qquad
m_6 = m_{10} = m_{12} = 1 \ (\text{and } m_{15}=1,
\text{parity-excluded in } I_h).
\label{eq:mult}
\end{equation}
The worked table (each entry the five-term sum divided by 60; the
verify script computes every line):
\begin{center}
\begin{tabular}{c|cccccc|c|ccc}
$l$ & $E$ & $12C_5$ & $12C_5^2$ & $20C_3$ & $15C_2$ & sum/60 & $m_l$ &&& \\
\hline
1 & 3 & $12\varphi$ & $-12/\varphi$ & 0 & $-15$ & 0 & 0 &&& \\
2 & 5 & 0 & 0 & $-20$ & 15 & 0 & 0 &&& \\
3 & 7 & $-12/\varphi$ & $12\varphi$ & 20 & $-15$ & 0 & 0 &&& \\
4 & 9 & $-12\varphi$ & $12/\varphi$ & 0 & 15 & 0 & 0 &&& \\
5 & 11 & $-12$ & $-12$ & $-20$ & $-15$ & 0 & 0 &&& \\
6 & 13 & $12/\varphi$ & $-12\varphi$ & 20 & 15 & 60 & \textbf{1} &&& \\
\end{tabular}
\end{center}
(with $\varphi$ the golden ratio; $\chi_l(72^\circ) \in \{ \varphi,
1/\varphi, -1/\varphi, -\varphi, -1, \ldots\}$ per the character
formula).  Since the full symmetry of the cell is $I_h = I \times \{1,
-1\}$ and the boundary function is even, odd-$l$ content vanishes
identically.
\textbf{Consequence:} any $I_h$-invariant function on the sphere --- in
particular the cell's boundary radius $R(\omega)$ and every mode of the
trivial irrep --- has angular content only in $l \in \{0, 6, 10, 12,
16, \ldots\}$.  There is no icosahedrally-invariant quadrupole,
octupole, or anything below $l = 6$: the forbidden channels are
forbidden \emph{exactly}, by group theory, not smallness.

\subsection{The actual cell realizes the lemma}
\label{sec:content}

Projecting the dodecahedron's support function
$R(\omega) = \rho_{\rm in}/\max_f(\omega\cdot\hat n_f)$ onto real
spherical harmonics (Gauss--Legendre $\times$ uniform-$\varphi$
quadrature, $240\times480$): the forbidden channels $l = 1\ldots5, 7,
8, 9, 11$ sit at the quadrature floor ($\le 5\times10^{-6}$), and the
anisotropy is
\begin{equation}
\epsilon_6 = 5.10\times10^{-2},\qquad
\epsilon_{10} = 1.13\times10^{-2},\qquad
\epsilon_{12} = 1.20\times10^{-2},
\label{eq:eps}
\end{equation}
with mean radius $\langle R\rangle = 1.0951\,\rho_{\rm in}$.  The
group-theoretic prediction \eqref{eq:mult} is realized channel by
channel in the actual Voronoi cell.

\subsection{Selection Theorem}
\label{sec:selection}

\begin{theorem}[Icosahedral selection]
On the true Voronoi cell with the resonator boundary conditions of
Spin~II (node at the center, free end at the boundary), the eigenmodes
classify by irreps of $I_h$.  Within the \emph{trivial} irrep --- the
sector of the ZBW cloud's scalar breathing modes --- the spectrum
begins with the deformed radial Modes 1 and 2: the first
anisotropy-carrying invariant channel is $l = 6$, whose lowest radial
eigenvalue under the same boundary conditions is
\begin{equation}
k_{6,1}R = 8.2108 \;>\; \tfrac{3\pi}{2} = 4.712 = k_2 R,
\end{equation}
a $74\%$ margin (exact route: the relevant zero of $[x\,j_6(x)]'$,
$8.21084$; the paper's FD instrument gives $8.2112$, $+0.005\%$; a
third independent discretization in review returned $8.2122$ --- three
routes, one ordering).  Hence the second invariant mode of the true cell
\emph{is} the anchoring mode of Spin~II's minimum-energy lemma: the
lowest invariant mode with both an interior antinode and an interior
node.  The lattice must choose Mode 2.
\end{theorem}

\noindent\emph{Where the low-$l$ modes went.}  Globally (ignoring
symmetry), the modes with $l = 1$ ($k_{1,1}R = 2.744$, three states)
and $l = 2$ ($k_{2,1}R = 3.870$, five states) do lie between the radial
Modes 1 and 2 --- eight states, which is exactly why the
Section~\ref{sec:measurement} run found its radial modes at global
indices 1 and 10.  The Lemma of Section~\ref{sec:characters} is what
excludes them: no $l = 1\ldots5$ content can appear in any invariant
mode, so the interlopers live entirely in non-trivial irreps and cannot
serve as the cloud's scalar breathing state.  Symmetry does not remove
them from the spectrum; it removes them from the \emph{sector}.

\subsection{Protection Theorem}
\label{sec:protection}

\begin{theorem}[Icosahedral protection]
Write the boundary as $R(\omega) = \langle R\rangle(1 +
\sum_{l\ge6}\delta_l(\omega))$ per \eqref{eq:mult}.  For a radial
unperturbed mode, the first-order (Hadamard) eigenvalue shift is an
integral of an angularly \emph{constant} density against $\delta R$,
so it picks out only the $l = 0$ mean --- explicitly, every $l \ge 6$
component integrates to zero against the constant density because
$\int Y_l^m\, d\Omega = 0$ for $l \ge 1$ --- a pure scale change that
moves no fractional position.  The first-order eigenfunction
correction lies entirely in the $l \ge 6$ channels, which vanish under
the angular average; the angular-averaged profile, whose zeros and
extrema define the node and antinode, is therefore unchanged at first
order.  Fractional node and antinode positions, and the
mean-normalized eigenvalue, shift only at $O(\epsilon_6^2)$.
\end{theorem}

\noindent\emph{Quantitative closure with the measurement.}  The
theorem predicts residuals of size $\epsilon_6^2 = 2.6\times10^{-3}$.
Section~\ref{sec:measurement} measured a mean-normalized eigenvalue
residual of $3.1\times10^{-3}$ --- within $20\%$ of the prediction ---
and a node-position shift of $5\times10^{-4}$, well inside the
$\epsilon_6^2$ class.  (Stated as consistency, not decomposition: the
instrument's own declared approximation enters at the same order.  In
the panel-adopted stricter wording: the Protection Theorem provides an
order-of-magnitude consistency check; a non-perturbative bound is
still required for full protection --- its registered discharge path is
to propagate the instrument's declared approximations through the
perturbative calculation and show the resulting uncertainty envelope
contains the measured residuals.)
The two legs of \texttt{OPEN-QM-8} now cross-confirm: the symmetry
argument predicts the smallness the instrument measured, and the
instrument measures the smallness the symmetry argument predicts.

\subsection{Scope of the proof}

Honest boundaries: the Protection Theorem is perturbative in
$\epsilon_6$ (first-order exclusion exact by symmetry; the $O(\epsilon_6^2)$
statement is standard boundary perturbation theory).  The Selection
Theorem's gap comparison uses the spherical $l = 6$ eigenvalue as the
reference for the invariant anisotropy channel; at the cell's measured
$\epsilon_6 = 0.051$ with a $74\%$ gap margin, no realistic
second-order shift reorders the sector, and the non-perturbative
measurement of Section~\ref{sec:measurement} independently confirms
the ordering (three radial candidates, isotropy scores
$0.999$--$1.000$, exactly where the theorem puts them).  What remains
registerable is a fully non-perturbative bound --- a rigor upgrade,
not a gap in the physical argument; recorded as the residual note on
\texttt{OPEN-QM-8}.

% =====================================================================
\section{CPP physical mechanism and conventional mapping}
\label{sec:mechanism}

\textbf{What is physically happening.}  An unpaired $-e$CP moving
through the Dipole Sea captures a DP (Spin I): the inner $+$CP orbits at
$r_{\rm in}$, the outer $-$CP at $r_{\rm out}$.  The ZBW polarization
cloud around the bare CP is a compression pattern of the Sea whose
scalar (breathing) degree of freedom obeys the radial wave equation of
Spin II: the CP Exclusion Rule pins a node at the center, and the
cloud's thermal boundary at $r_{\rm th}$ is a free end.  The captured
DP anchors to the standing wave's structure --- $+$CP at the interior
antinode, $-$CP at the interior node --- which is what makes Mode 2,
the lowest mode possessing both features, the physically selected
state, and what fixes $r_{\rm out} = 2 r_{\rm in}$ and hence
$L = \hbar/2$.  This paper's contribution: the wave does not live in a
continuum but inside one Grid Point's territory --- the regular
dodecahedral Voronoi cell of the 600-cell lattice --- and the cell's
icosahedral symmetry \emph{forbids} every non-radial competitor below
the $l = 6$ channel, forcing the selection the continuum papers
assumed.

\textbf{How this maps to conventional physics.}  The correspondence is
structural, not literal:

\begin{center}
\begin{tabular}{p{0.44\textwidth}|p{0.46\textwidth}}
\textbf{CPP element} & \textbf{Conventional counterpart} \\
\hline
ZBW cloud scalar breathing mode & Zitterbewegung of the Dirac electron
(the $\sim 2mc^2/\hbar$ internal oscillation) \\
Captured-DP orbital angular momentum $L = \hbar/2$ & Electron spin
$s = \tfrac12$ (intrinsic, not orbital, in the conventional account) \\
CP Exclusion node at $r = 0$ & Boundary condition selecting the mode
family (no conventional analogue --- axiom-level) \\
Thermal boundary free end at $r_{\rm th} = \hbar/2m_ec$ & Reduced
Compton wavelength scale of the electron \\
Trivial irrep of $I_h$ (invariant sector) & $s$-wave ($l = 0$) sector of
a central-potential problem \\
Icosahedral forbidden channels $l = 1..5$ & Crystal-field selection
rules (here: exact, from character theory) \\
$O(\epsilon_6^2)$ position protection & Second-order perturbation
theory under a symmetry-restricted perturbation \\
\end{tabular}
\end{center}

The conventional account postulates spin-$\tfrac12$ as an irreducible
quantum number; CPP derives the same magnitude from a resonator whose
geometry the lattice enforces.  The discriminating content is not the
value (both give $\hbar/2$) but the mechanism's collateral structure:
the mode spectrum, the scale bridge to Spin I, and the icosahedral
selection rules of Section~\ref{sec:analytic}.

% =====================================================================
\section{Epistemic ledger}
\label{sec:ledger}

\begin{itemize}[leftmargin=2em]
\item \textbf{Exact:} the resonator spectrum \eqref{eq:spectrum} and
Mode-2 ratio (Spin~II, restated and re-verified here two ways); the
scale-bridge identity \eqref{eq:bridge}; the discreteness bound.
\item \textbf{Established numerically:} FD eigenvalues to $10^{-4}$
across eight modes; Mode-2 geometry to $10^{-3}$; committed CSV
regenerated bit-identically.
\item \textbf{Measured (V0.1):} the lattice-cell selection question,
numerical leg --- \textsc{mode2-recovered} on the true Voronoi cell at
both densities under the pre-frozen verdict rule
(Section~\ref{sec:measurement}); the committed March instrument's
closed-problem defect stands diagnosed (Section~\ref{sec:instrument}).
\item \textbf{Proved (V0.2):} the analytic leg --- exact invariant
content ($m_l = 0$, $1\le l\le5$; first channel $l=6$), the Selection
Theorem ($k_{6,1}R = 8.211 > 3\pi/2$: the second invariant mode is
Mode 2), and the Protection Theorem ($O(\epsilon_6^2)$ positions;
predicted $0.0026$ vs measured $0.0031$).  Residual (rigor upgrade
only): a fully non-perturbative bound.
\item \textbf{Ruled:} A1, the Voronoi domain --- true cell (regular
dodecahedron); 24-cell retired as a prior worker's initiative
(Section~\ref{sec:domain}).
\item \textbf{Inherited:} everything Spin~I and Spin~II inherit; in
particular the $L = \hbar/2$ chain is complete at the continuum level
(Spin~I $+$ Spin~II) independently of this paper's open capstone.
\end{itemize}

% =====================================================================
\section*{Provenance}

Drafted at founder request (Session 149) from the arc's own materials:
the development-notes dialogue (Thomas Lee Abshier with Claude and
Grok, March 2026), the committed eigenvalue script and its data and
figure.  No prior Spin~III paper existed (full-history finding,
Patch 3232 record).  Verify:
\texttt{scripts/3234\_verify\_spin3\_v0.py} --- 10/10, including the
reproducible Part-3 diagnosis.

\begin{thebibliography}{9}
\bibitem{spin1} T.~L. Abshier, \emph{Emergent Spin-$\tfrac12$ from
Captured Dipole Particle Orbital Geometry} (spin sub-arc, SPIN-1; in-arc name Spin~I).
\bibitem{spin2} T.~L. Abshier, \emph{Derivation of the Standing-Wave
Sub-Harmonic Condition for Captured Dipole Particle Orbits} (spin
sub-arc, SPIN-2; in-arc name Spin~II).
\bibitem{zbwmass} T.~L. Abshier, \emph{ZBW $\hbar$ and Mass Units}
(SR companion c04) --- source of the thermal-boundary condition.
\end{thebibliography}

\end{document}
